% !TEX root = ../main.tex
\chapter{Recurrent Networks}
\label{ch:rnn_family}

Recurrent neural networks process variable-length sequences by
applying the same transformation at each timestep, threading a
\emph{hidden state} that carries information across the sequence.
Before the Transformer's dominance, recurrent networks (RNN, LSTM,
GRU) were the standard architecture for sequence modeling: machine
translation, speech recognition, language modeling, time-series
forecasting. They retain a niche in modern systems where causal
streaming or low memory matter (online speech recognition, audio
codecs, certain reinforcement-learning policies).

\Lucid\ exposes the standard family at
\code{lucid.nn.RNN}, \code{lucid.nn.LSTM}, \code{lucid.nn.GRU},
mirroring \refframework. This chapter describes the
mathematical recurrence of each, the implementation, the
\code{proj\_size} extension, the bidirectional wrapper, and the
packing utilities for variable-length sequences.

\section{The Generic Recurrence}
\label{sec:rnn_family:generic}

All RNN variants share the form
\[
h_t \;=\; f(x_t, h_{t-1}; \theta),
\]
where $x_t$ is the input at timestep $t$, $h_t$ is the hidden
state, and $f$ is a parameterised transformation. The differences
between RNN/LSTM/GRU are in the choice of $f$.

The output at each timestep is some function of $h_t$ (often
$h_t$ itself or a linear projection thereof). The full sequence
of outputs $(y_1, y_2, \ldots, y_T)$ is what the layer returns,
along with the final hidden state $h_T$ that the user can pass to
the next layer or the next forward call.

\section{Vanilla RNN}
\label{sec:rnn_family:vanilla_rnn}

The simplest recurrence:
\[
h_t \;=\; \tanh(W_{ih}\, x_t + b_{ih} + W_{hh}\, h_{t-1} + b_{hh}).
\]

\code{nn.RNN(input\_size, hidden\_size)} constructs this layer.
For multi-layer stacks (a sequence of RNN layers each consuming
the previous one's output), \code{num\_layers > 1}.

\subsection{The vanishing/exploding gradient}
\label{ssec:rnn_family:vanishing}

The vanilla RNN's gradient through time involves $T$ products of
$W_{hh}$ Jacobians, which either vanishes or explodes
exponentially in $T$ unless the spectral radius of $W_{hh}$ is
exactly 1. This is the original motivation for LSTM and GRU,
both of which introduce gating to mitigate the issue.

\subsection{Back-propagation through time (BPTT)}
\label{ssec:rnn_family:bptt}

The training of a recurrent network on a sequence of length $T$
involves unrolling the recurrence and applying standard reverse-
mode AD over the unrolled DAG. Concretely, the loss
$\mathcal{L} = \sum_{t=1}^T \ell_t(h_t)$ has gradient with
respect to $h_{t-1}$:
\[
  \frac{\partial \mathcal{L}}{\partial h_{t-1}}
  \;=\;
  \frac{\partial h_t}{\partial h_{t-1}} \cdot
  \frac{\partial \mathcal{L}}{\partial h_t}
  + \frac{\partial \ell_{t-1}}{\partial h_{t-1}}.
\]
Iterating this gives a chain whose Jacobian product accumulates
across the entire sequence. For vanilla RNN with
$h_t = \tanh(W_{hh} h_{t-1} + W_{ih} x_t + b)$:
\[
  \frac{\partial h_t}{\partial h_{t-1}}
  \;=\;
  \mathrm{diag}(1 - h_t^2) \cdot W_{hh}.
\]
The spectral radius $\rho(W_{hh})$ governs the long-term gradient
behaviour: $\rho < 1$ shrinks gradients to zero (vanishing),
$\rho > 1$ blows them up (exploding). The tanh derivative
$1 - h_t^2 \leq 1$ only worsens the vanishing side.

\subsection{Truncated BPTT}
\label{ssec:rnn_family:truncated_bptt}

For very long sequences, full BPTT is memory-prohibitive (the
forward activations of all $T$ timesteps must be stored). The
standard workaround: \emph{truncated BPTT} (TBPTT) backpropagates
only $k$ steps, treating $h_{t-k}$ as a fresh input (no gradient
flow further back).

\begin{lstlisting}[style=lucid,language=LucidPython]
hidden = None
for chunk in sequence.split(chunk_size=64):  # 64-step chunks
    out, hidden = lstm(chunk, hidden)
    loss = criterion(out, targets[chunk])
    loss.backward()
    optimizer.step()
    hidden = tuple(h.detach() for h in hidden)  # break gradient
\end{lstlisting}

The \code{.detach()} call is the truncation: it severs the
autograd graph between chunks. The hidden state itself
continues; only the gradient stops.

TBPTT introduces a bias (gradient is computed for an in-truth
finite window) in exchange for tractable memory; the bias is
usually acceptable.

\section{LSTM}
\label{sec:rnn_family:lstm}

The Long Short-Term Memory cell maintains two states: a hidden
state $h_t$ (the output) and a cell state $c_t$ (the internal
memory). Four gates control information flow:
\begin{align*}
i_t &\;=\; \sigma(W_{ii}\, x_t + W_{hi}\, h_{t-1} + b_i)
   \quad\text{(input gate)} \\
f_t &\;=\; \sigma(W_{if}\, x_t + W_{hf}\, h_{t-1} + b_f)
   \quad\text{(forget gate)} \\
g_t &\;=\; \tanh(W_{ig}\, x_t + W_{hg}\, h_{t-1} + b_g)
   \quad\text{(candidate)} \\
o_t &\;=\; \sigma(W_{io}\, x_t + W_{ho}\, h_{t-1} + b_o)
   \quad\text{(output gate)} \\
c_t &\;=\; f_t \odot c_{t-1} + i_t \odot g_t \\
h_t &\;=\; o_t \odot \tanh(c_t).
\end{align*}

The gating equations are the essence: the cell state $c_t$ evolves
multiplicatively (controlled by $f_t$ and $i_t$), which allows
gradients to flow back through many timesteps without
vanishing/exploding under typical parameter scales.

\subsection{The Lucid API}
\label{ssec:rnn_family:lstm_api}

\begin{lstlisting}[style=lucid,language=LucidPython]
lstm = nn.LSTM(input_size=128, hidden_size=256, num_layers=2,
               batch_first=True, dropout=0.1)
# x.shape = (B, T, 128)
output, (h_n, c_n) = lstm(x)
# output.shape = (B, T, 256)  -- per-timestep hidden states
# h_n.shape = (num_layers, B, 256)
# c_n.shape = (num_layers, B, 256)
\end{lstlisting}

The \code{batch\_first=True} convention puts batch first in input
and output shapes (the \refframework\ default is
\code{batch\_first=False}, i.e., $(T, B, \cdot)$; we expose both
but recommend \code{True} as more intuitive).

\subsection{The proj\_size extension}
\label{ssec:rnn_family:proj_size}

A common trick for very wide LSTMs: project the hidden state to a
lower-dimensional output. \code{nn.LSTM(..., proj\_size=p)} adds
a learnable projection $W_{hr} h_t$ that reduces the output from
\code{hidden\_size} to $p$:
\[
h_t \;=\; W_{hr}\,(o_t \odot \tanh(c_t)),
\quad h_t \in \mathbb{R}^p.
\]

The $c_t$ remains full-size; only the output side is projected.
Used in the original Tacotron and in some speech-recognition
backbones to reduce parameter count without sacrificing memory
capacity.

\subsection{Bidirectional}
\label{ssec:rnn_family:bidirectional}

\code{nn.LSTM(..., bidirectional=True)} runs two LSTMs, one
forward and one backward, and concatenates their outputs. The
output dimension doubles. Used for non-causal tasks (named entity
recognition, sequence classification) where the future context is
available.

Implementation: the layer holds two complete parameter sets
(suffixed \code{\_reverse}), runs the input through both
directions in parallel, and concatenates the per-timestep
outputs along the feature axis. The hidden states
$(h_n, c_n)$ stack along a new leading axis:
\code{(num\_layers $\times$ 2, B, hidden\_size)}.

Bidirectional RNNs are incompatible with causal streaming (the
backward direction needs the entire sequence). In
encoder-decoder architectures, the encoder is typically
bidirectional and the decoder unidirectional --- the decoder
generates timestep by timestep, while the encoder has full
input context.

\subsection{Layer-wise normalisation in LSTM}
\label{ssec:rnn_family:lstm_ln}

Long sequences can drift activations into saturation; layer
normalisation inside the LSTM cell mitigates this. The
modified equations apply LayerNorm to the four pre-activation
linear projections:
\begin{align*}
  \tilde{i}_t &= \mathrm{LN}_i(W_{ii} x_t + W_{hi} h_{t-1}), \\
  i_t &= \sigma(\tilde{i}_t),
\end{align*}
and similarly for $f, g, o$. The cell-state update applies
LayerNorm to $c_t$ as well.

\Lucid\ does not currently provide a fused LayerNorm-LSTM
module; users implement it as a custom \code{nn.Module}
subclass when needed. The standard \code{nn.LSTM} is the
unnormalised form.

\section{GRU}
\label{sec:rnn_family:gru}

The Gated Recurrent Unit is a simpler alternative to LSTM:
\begin{align*}
r_t &\;=\; \sigma(W_{ir}\, x_t + W_{hr}\, h_{t-1} + b_r)
   \quad\text{(reset gate)} \\
z_t &\;=\; \sigma(W_{iz}\, x_t + W_{hz}\, h_{t-1} + b_z)
   \quad\text{(update gate)} \\
n_t &\;=\; \tanh(W_{in}\, x_t + r_t \odot (W_{hn}\, h_{t-1}) + b_n) \\
h_t &\;=\; (1 - z_t) \odot n_t + z_t \odot h_{t-1}.
\end{align*}

GRU has three gates (vs LSTM's four) and one state (vs LSTM's
two). It has fewer parameters and is faster to train; empirical
results on most benchmarks are comparable to LSTM. The choice
between them is application-dependent.

\section{The Hidden-State Convention}
\label{sec:rnn_family:hidden_convention}

The initial hidden state $h_0$ (and cell state $c_0$ for LSTM) is
zero by default. The user can pass a custom initial state:

\begin{lstlisting}[style=lucid,language=LucidPython]
h_0 = lucid.zeros(num_layers, batch_size, hidden_size)
c_0 = lucid.zeros(num_layers, batch_size, hidden_size)
output, (h_n, c_n) = lstm(x, (h_0, c_0))
\end{lstlisting}

A common use case: continuing a sequence across multiple forward
calls (truncated backpropagation through time, online inference).
Pass the previous call's $(h_n, c_n)$ as the next call's initial
state, with \code{.detach()} on each to break the gradient.

\section{Variable-Length Sequences}
\label{sec:rnn_family:variable_length}

A batch of sequences may have different lengths. The two standard
approaches:

\begin{itemize}
  \item \textbf{Padding}: pad shorter sequences to the longest
    with zeros, run the RNN over the padded tensor. Wasteful
    compute on padding positions; the user must mask out the
    padding in any loss computation.
  \item \textbf{Packing}: use \code{PackedSequence} to compact
    the variable-length batch into a single tensor with length
    metadata; the RNN respects the lengths.
\end{itemize}

\subsection{PackedSequence}
\label{ssec:rnn_family:packed_sequence}

\code{lucid.nn.utils.rnn.pack\_padded\_sequence(padded, lengths)}
converts a padded batch into a \code{PackedSequence}. The RNN
operates on the packed form and produces another packed sequence;
\code{pad\_packed\_sequence(packed)} unpacks back.

\begin{lstlisting}[style=lucid,language=LucidPython]
from lucid.nn.utils.rnn import pack_padded_sequence, pad_packed_sequence

# Sort by length descending (required by packing API)
lengths, perm = lengths.sort(descending=True)
x = x[perm]

packed = pack_padded_sequence(x, lengths, batch_first=True)
output_packed, (h, c) = lstm(packed)
output, _ = pad_packed_sequence(output_packed, batch_first=True)
\end{lstlisting}

The packed form avoids wasted compute on padding positions; the
speedup is proportional to the variance in sequence length.

\section{Implementation Notes}
\label{sec:rnn_family:impl_notes}

The RNN family is implemented in \code{lucid/\_C/nn/LSTM.cpp} (and
analogous files for RNN and GRU). The forward unrolls the
recurrence in C++ and dispatches per-timestep through the matrix
multiplication backend.

\subsection{The cuDNN-equivalent fast path}
\label{ssec:rnn_family:fast_path}

\refframework's \code{nn.LSTM} on CUDA dispatches through cuDNN's
fused multi-step kernels, which are substantially faster than
naive per-timestep unrolling. \Lucid\ on Apple silicon does not
have an equivalent fused kernel; the implementation is the naive
unrolled form.

The performance gap is moderate (2--4$\times$ slower than cuDNN's
fused LSTM for sequence lengths $> 50$). For most production
workloads on \Lucid, this is acceptable because the modern
alternative is the Transformer (\chref{ch:attention}), which is
faster and produces better results on most language tasks.

\subsection{The flatten\_parameters call}
\label{ssec:rnn_family:flatten_parameters}

The \refframework\ idiom \code{model.flatten\_parameters()} is a
no-op in \Lucid; it exists on the API for compatibility but does
nothing. (\refframework's version optimises cuDNN's memory
access; without cuDNN, the optimisation is irrelevant.)

\section{Numerical Considerations}
\label{sec:rnn_family:numerical}

\subsection{Saturating activations}
\label{ssec:rnn_family:saturating}

LSTM and GRU use sigmoid and tanh activations, both of which
saturate at large magnitudes (gradient $\to 0$). For very long
sequences, the gating equations can drive intermediate values
into saturation, causing the gradient to vanish anyway.

The standard mitigation is gradient clipping (\code{clip\_grad\_norm}
in \code{lucid.nn.utils}; \chref{ch:optimizers}), which caps the
backward pass's gradient magnitude.

\subsection{Initialisation}
\label{ssec:rnn_family:rnn_init}

LSTM weights are initialised from
$\mathcal{U}(-1/\sqrt{H}, 1/\sqrt{H})$ for hidden size $H$. The
forget gate bias is often initialised to 1 (rather than 0) to
encourage the model to remember at the start of training; this is
the \emph{Jozefowicz initialisation}. \Lucid\ does not apply this
by default but the user can:

\begin{lstlisting}[style=lucid,language=LucidPython]
for name, param in lstm.named_parameters():
    if 'bias_hh' in name:
        H = param.shape[0] // 4
        param.data[H:2*H].fill_(1.0)  # forget gate
\end{lstlisting}

\section{Performance Notes}
\label{sec:rnn_family:perf}

Measured throughput on M3~Max for a batch-size-32, sequence-length-
256 LSTM with hidden size 512:

\begin{itemize}
  \item Forward: 18~ms / step.
  \item Backward: 42~ms / step.
  \item Total: $\sim$60~ms / step.
\end{itemize}

For comparison, a Transformer layer of the same effective
capacity runs in $\sim$8~ms forward / $\sim$22~ms backward on
the same hardware. The LSTM's 3$\times$ slowdown comes from the
sequential dependency (each timestep waits for the previous's
hidden state), which prevents parallelisation across timesteps.

\subsection{Where the time goes}
\label{ssec:rnn_family:perf_breakdown}

A per-timestep LSTM cell does:
\begin{itemize}
  \item 8 matrix multiplications ($4 \times W_{ix} x + 4 \times W_{ih} h$),
    one per gate plus the candidate.
  \item 4 sigmoids and 1 tanh.
  \item 3 elementwise products (the gating).
  \item 1 cell-state update and 1 hidden-state computation.
\end{itemize}
Total: $\sim$\,20 op invocations per timestep, each with the
per-op dispatch overhead. For 256 timesteps, that's
$\sim$\,5000 op calls per LSTM forward --- the dispatch
overhead alone is several milliseconds.

The compile path (\chref{ch:compile_design}) eliminates most of
this: the entire unrolled LSTM becomes one MPSGraph
submission. Compile speedup on the same configuration:
60 ms $\to$ 28 ms ($2.1\times$).

\subsection{Sequence-length scaling}
\label{ssec:rnn_family:seq_scaling}

\begin{table}[t]
\centering
\small
\begin{tabular}{lrr}
\toprule
Seq length & LSTM (ms) & Transformer (ms) \\
\midrule
64    & 17  & 12 \\
128   & 32  & 16 \\
256   & 60  & 30 \\
512   & 120 & 88 \\
1024  & 245 & 310 (attention dominates) \\
2048  & 490 & 1180 \\
\bottomrule
\end{tabular}
\caption[Sequence-length scaling: LSTM vs Transformer.]{Per-step
training time on M3~Max for an LSTM and a single
Transformer encoder block of equivalent hidden size as
sequence length grows. LSTM scales linearly with $T$;
Transformer scales as $T^2$ in attention. The crossover is
around $T = 512$ on M3 Max; beyond that, RNNs are competitive
again.}
\label{tab:rnn_family:seq_scaling}
\end{table}

\section{Comparison of the Three}
\label{sec:rnn_family:comparison}

\begin{table}[t]
\centering
\small
\begin{tabular}{lccc}
\toprule
Property & RNN & LSTM & GRU \\
\midrule
States per cell  & $h$ (1) & $h, c$ (2) & $h$ (1) \\
Gates per cell   & 0       & 4          & 3      \\
Params (per cell) & $H(H+D)$ & $4H(H+D)$ & $3H(H+D)$ \\
Forward FLOPs/step & $2H(H+D)$ & $8H(H+D)$ & $6H(H+D)$ \\
Vanishing-gradient & strong  & mild      & mild   \\
Recommended use & demos / tiny & general & general \\
\bottomrule
\end{tabular}
\caption[RNN/LSTM/GRU comparison.]{Parameter and FLOP counts per
recurrent cell, with hidden size $H$ and input size $D$. LSTM
has roughly 1.3$\times$ the parameters and FLOPs of GRU; in
practice the two perform comparably on most tasks. Vanilla RNN
is rarely used in production because of vanishing-gradient
issues.}
\label{tab:rnn_family:comparison}
\end{table}

\section{Implementation: The Per-Timestep Loop}
\label{sec:rnn_family:per_timestep}

A sketch of the unrolled-loop implementation that
\code{lucid.nn.LSTM} executes:

\begin{lstlisting}[style=lucid,language=LucidPython]
def lstm_forward(input, h_0, c_0, w_ih, w_hh, b_ih, b_hh):
    T = input.size(0)
    H = h_0.size(-1)
    h, c = h_0, c_0
    outputs = []
    for t in range(T):
        x_t = input[t]
        # 8 matmuls collapsed into 2 (gates concatenated)
        gates = (x_t @ w_ih.T + b_ih) + (h @ w_hh.T + b_hh)
        # gates shape: (B, 4H); split into i, f, g, o
        i, f, g, o = gates.chunk(4, dim=-1)
        i = i.sigmoid()
        f = f.sigmoid()
        g = g.tanh()
        o = o.sigmoid()
        c = f * c + i * g
        h = o * c.tanh()
        outputs.append(h)
    return lucid.stack(outputs, dim=0), (h, c)
\end{lstlisting}

The trick \code{chunk(4)} concatenates the four gates'
matrix multiplications into one big GEMM, which is more
efficient than four separate GEMMs. Both the input-to-hidden
projection and the hidden-to-hidden projection get this
treatment.

\subsection{Multi-layer stacks}
\label{ssec:rnn_family:multi_layer}

For \code{num\_layers > 1}, the layers are applied sequentially:
the output of layer $\ell$ at every timestep becomes the input
of layer $\ell + 1$. Each layer has its own parameters.

The naive implementation has a dependency: layer 2 timestep 1
needs layer 1 timestep 1. But layer 2 timestep 2 only needs
layer 1 timestep 1 and 2 to have completed --- it does not
need layer 2 timestep 1 to have finished (only its $h_1$ is
needed). The intra-layer dependency is per-layer; the
inter-layer dependency is feed-forward.

This gives a parallelism opportunity: in principle, the
layer-2 timesteps can run in parallel with the layer-1
timesteps, with a small lag. Apple's MPS does not currently
exploit this; the implementation is sequential per layer per
timestep.

\subsection{Bidirectional implementation}
\label{ssec:rnn_family:bidir_impl}

The two directions of a bidirectional LSTM run as two
independent unrolled loops; the outputs are concatenated along
the feature axis at each timestep. The backward direction
reverses the time axis:

\begin{lstlisting}[style=lucid,language=LucidPython]
out_fwd, (h_fwd, c_fwd) = lstm_forward(input, h_0_fwd, c_0_fwd, ...)
out_bwd, (h_bwd, c_bwd) = lstm_forward(input.flip(0), h_0_bwd, c_0_bwd, ...)
out_bwd = out_bwd.flip(0)
output = lucid.cat([out_fwd, out_bwd], dim=-1)
\end{lstlisting}

The two directions can in principle run in parallel; the
current implementation runs them sequentially.

\section{When to Use RNN/LSTM/GRU}
\label{sec:rnn_family:when_to_use}

Despite the Transformer's dominance, recurrent networks remain
the right choice in three scenarios:

\begin{enumerate}
  \item \textbf{Causal streaming}: a model that must produce
    output incrementally as input arrives (online speech
    recognition, audio codec, real-time captioning) needs the
    O(1) per-step state that RNNs provide. Transformers' KV cache
    is an alternative but has its own complexity.
  \item \textbf{Very long sequences}: Transformers scale as
    $O(T^2)$ in attention; RNNs scale as $O(T)$. For sequences
    beyond $\sim 10^4$ tokens, the asymptotic difference matters.
  \item \textbf{Constrained inference memory}: an LSTM has a
    fixed-size hidden state regardless of sequence length;
    Transformers' KV cache grows with sequence length. For
    deployment with strict memory budgets, the constant is
    valuable.
\end{enumerate}

For most other settings (machine translation, summarisation,
classification on bounded-length input), Transformers
(\chref{ch:attention}) are the modern choice.

\section{Summary and Forward References}
\label{sec:rnn_family:summary}

\Lucid's recurrent network family (\code{nn.RNN}, \code{nn.LSTM},
\code{nn.GRU}) implements the standard gated recurrences with
multi-layer stacking, bidirectional wrapping, LSTM projection,
and variable-length sequence packing. The implementation lacks
the cuDNN-equivalent fused fast path; performance is a few
$\times$ slower than \refframework\ on CUDA but acceptable for
typical workloads.

The next chapter, \chref{ch:attention}, treats the attention
mechanism that has largely displaced RNNs in modern NLP.
\chref{ch:losses} treats loss functions. \chref{ch:initialization}
treats weight initialisation theory.

\begin{lucidnote}
For the user: the RNN family API is drop-in compatible with
\refframework. The semantics of \code{(h\_n, c\_n)} returns, the
\code{batch\_first} convention, \code{bidirectional},
\code{proj\_size}, and the packing utilities all work identically.
\end{lucidnote}
